% =========================
% puzzles/p01.tex
% Final — High-Card Poker: "Lower Than Six"
% =========================

% Use the PuzzleChapterIntro macro to supply an introduction in the parchment box.
\PuzzleChapterIntro{High-Card Poker: ``Lower Than Six''}{The Final Wager}{This is a one-hand inference puzzle: a single spoken claim about a hidden card is a noisy signal, and you must update the hidden-card distribution before you can compute the chance of a strict win. It tests clean sample-space counting in a 52-card deck, Bayes conditioning under an explicit truth/lie rate, and careful treatment of ``win'' versus ``tie.'' The only required tools are elementary counting and conditional probability.}

\Lore{The Grand Archive Casino hosts a simple game: the higher card wins, and ties pay nothing.

Each table is sleeved in green felt, inlaid with walnut dark as old tea and a seam of obsidian that catches the light like a blade. The cards are linen-backed, worn smooth by centuries of hands. The chandeliers do not merely flicker; they \emph{remember}. Here, a remark is never just a remark. It is a move.

Around the chamber hang paintings---twenty of them---none newer than the truths they suggest. The house permits ornament, but it forbids invention. Every detail in this room is placed on purpose, and every placement is recorded.

\begin{quote}
\textit{``A faithful replica of Van Gogh's \emph{Cafe Terrace at Night}. Van Gogh portrayed the reality of life through his painting---here, all descriptions are literal.''}
\end{quote}

To \textbf{Cyrene}'s right, the cafe glows under an indigo sky. Three figures sit at a table near the center, faces unfinished, conversation lost to time. A glass catches a thin blade of gold. The frame is heavy gilt---polished, too perfect, and therefore not old. It is expensive in the way a mask is expensive.

To her left hangs something that does not bother with beauty. It is a panel, pigment laid straight onto rough grain, held in a cradle of dark wood. No signature. No gallery stamp. A circular table, mid-shuffle. No hands. An overturned glass. The shadows fall in quiet blocks, the kind that look, at a distance, like numbers.

At the bottom of that frame, etched into a narrow strip of brass---dulled by touch---is a line the house has not polished away:

\begin{quote}
\textit{Did you know? In over 10{,}000 observed games, comparative statements spoken \emph{during play} are correct only one time in four.}
\end{quote}

The dealer burns one card. The room goes still enough to hear the felt breathe.}

\Problem

\textbf{High-card rules.} Each player holds exactly one card from a standard 52-card deck. Ranks are ordered
\[
2<3<\cdots<10<J<Q<K<A.
\]
Suits do not matter. \textbf{Cyrene} \emph{wins} if and only if her rank is \emph{strictly higher} than \textbf{Irin}'s rank. (A tie is not a win.)

\medskip

Cyrene looks down at \(8\clubsuit\). Irin draws one card uniformly at random from the remaining \(51\) cards and keeps it hidden.

\medskip

Irin makes exactly one comparative claim about his rank relative to \(6\): either ``my card is lower than \(6\)'' or ``my card is \(6\) or higher.'' Cyrene hears Irin say:
\[
\text{``My card is lower than \(6\), Cyrene. Your move.''}
\]

\medskip

\textbf{Plaque rule.} Treat the brass-plaque statistic as binding for this exact two-option claim.

\VaultDirective{What is the probability that Cyrene wins given Irin's claim? Give your answer as a percentage rounded to the nearest whole percent.}

\SolutionPage

Let \(L\) be the event ``Irin's rank is below \(6\)'' (i.e.\ \(2,3,4,5\)), and let \(S\) be the event ``Irin says `lower than \(6\)'.''

By the plaque rule,
\[
P(S\mid L)=\frac14,\qquad P(S\mid \neg L)=\frac34.
\]

\medskip

\textbf{Priors.}

Cyrene  holds \(8\clubsuit\), so Irin draws uniformly from the remaining \(51\) cards. There are \(16\) cards of ranks \(2\)–\(5\), hence
\[
P(L)=\frac{16}{51},\qquad P(\neg L)=\frac{35}{51}.
\]

\medskip

\textbf{Bayes.}

\[
P(S)=\frac14\cdot\frac{16}{51}+\frac34\cdot\frac{35}{51}=\frac{121}{204}.
\]
So
\[
P(L\mid S)=\frac{(\frac14)(\frac{16}{51})}{\frac{121}{204}}=\frac{16}{121},
\qquad
P(\neg L\mid S)=\frac{105}{121}.
\]

\medskip

\textbf{Win chance.}

If \(L\) occurs, then Irin holds a \(2,3,4,\) or \(5\), so Cyrene's \(8\) wins surely:
\[
P(\text{win}\mid L,S)=1.
\]

If \(\neg L\) occurs, Irin's card lies among the \(35\) remaining cards with ranks \(6,7,8,9,10,J,Q,K,A\) (excluding \(8\clubsuit\)).
Cyrene's \(8\) beats only ranks \(6\) and \(7\), i.e.\ \(8\) winning cards, so
\[
P(\text{win}\mid \neg L,S)=\frac{8}{35}.
\]

Therefore,
\[
P(\text{win}\mid S)=\frac{16}{121}\cdot 1+\frac{105}{121}\cdot\frac{8}{35}
=\frac{40}{121}\approx 0.33058.
\]

\FinalAnswer{33\%}

\NextPage
